% 02_introduction.tex
\section{Introduction}
\label{sec:intro}

EU AI Act Article~50~\citep{euaiact2024} requires a machine-readable
mark on AI-generated text by August 2026. It does not say what
counts as marked: which false-positive rate, which text length,
which quality cost.
Existing schemes do not fill the gap. Logit-bias
watermarks~\citep{kirchenbauer2023watermark} tune a bias and a
green-list size by grid search.
Undetectability-frontier results~\citep{christ2024undetectable}
bound what any scheme can achieve, but yield no deployable
construction. Passive
detection~\citep{mitchell2023detectgpt,tian2023gptzero} has no
controllable frontier at all.

\begin{figure*}[t]
  \centering
  \begin{tikzpicture}[
    x=1cm, y=1cm,
    paneltitle/.style={font=\small\bfseries, anchor=base, inner sep=0pt,
                       text=black!85},
    paneltitlewest/.style={font=\small\bfseries, anchor=base west,
                           inner sep=0pt, text=black!85},
    axisline/.style={black!60, line width=0.4pt},
    gridline/.style={black!12, line width=0.3pt, dotted},
    pt/.style={circle, draw=black!55, line width=0.4pt,
               minimum size=2.2mm, inner sep=0pt},
    bigpt/.style={circle, draw=black!50, line width=0.45pt,
                  minimum size=4.5mm, inner sep=0pt},
    wmline/.style={line width=1.4pt, green!55!black,
                   line cap=round, line join=round},
    rndline/.style={line width=1.4pt, red!75!black, dashed,
                    line cap=round, line join=round},
    yticklbl/.style={font=\scriptsize, anchor=east, inner sep=1.5pt},
    xticklbl/.style={font=\scriptsize, anchor=north, inner sep=1.5pt},
    axislabel/.style={font=\footnotesize, color=black!75},
    score/.style={font=\normalsize\bfseries, anchor=west, inner sep=0pt},
    flowarr/.style={->, line width=0.9pt, black!75,
                    >={Stealth[length=2.6mm,width=2mm]}},
  ]
    \definecolor{c0}{HTML}{E15A5A}
    \definecolor{c1}{HTML}{F0A030}
    \definecolor{c2}{HTML}{E8C12F}
    \definecolor{c3}{HTML}{4FAE5C}
    \definecolor{c4}{HTML}{4A8FD3}

    % ===================================================================
    % Three-panel horizontal pipeline:
    %   A. clockwork (the hook, dominant visual)
    %   B. generation step (logits -> mask -> token)
    %   C. detection (watermarked vs random sequences + verdict)
    %
    % uniform header baseline at y = 4.95
    % uniform content vertical band roughly y in [0.4, 4.4]
    % ===================================================================

    % =========== PANEL A: clockwork ring (dominant hook) ===========
    %    column x in [0, 3.6], donut centred at (1.80, 2.55)
    \node[paneltitle] at (1.80, 4.95) {\large A.~Clockwork chain};
    \node[font=\footnotesize, color=black!75, anchor=base]
      at (1.80, 4.30)
      {$\sigma_k(t) = \mathrm{SHA\text{-}256}(k\,\Vert\,t)\bmod S$};

    \begin{scope}[shift={(1.80, 2.30)}]
      % ring sectors (outer 1.45, inner 0.65 -> wider colour band)
      \foreach \i [evaluate=\i as \stt using {90 - \i*72},
                   evaluate=\i as \fin using {90 - (\i+1)*72}] in {0,1,2,3,4} {
        \fill[c\i!55, draw=black!45, line width=0.45pt]
          (\stt:1.45) arc[start angle=\stt, end angle=\fin, radius=1.45]
          -- (\fin:0.65) arc[start angle=\fin, end angle=\stt, radius=0.65]
          -- cycle;
      }
      \draw[black!25, line width=0.3pt] (0,0) circle (1.45);
      \draw[black!25, line width=0.3pt] (0,0) circle (0.65);
      \foreach \i [evaluate=\i as \mid using {90 - (\i+0.5)*72}] in {0,1,2,3,4} {
        \node[font=\small\bfseries, color=black!85]
          at (\mid:1.05) {$\mathcal V_{\i}$};
      }
      % bold clockwise rotation indicator in the inner hole.
      % Sweeps 330 deg clockwise from upper-right (60 deg) down and
      % around to the top (90 deg in normal coords, -270 here). The
      % arrowhead therefore sits at the TOP of the circle and points
      % RIGHT -- the canonical "refresh / rotate clockwise" symbol.
      \draw[->, line width=2pt, black!85, line cap=round,
            >={Stealth[length=3.6mm,width=3mm,inset=0.4mm]}]
        (60:0.45) arc[start angle=60, end angle=-270, radius=0.45];
    \end{scope}

    % =========== PANEL B: one generation step (vertical flow) ===========
    %    column x in [4.0, 8.2], centre at x = 6.10
    %    vertical flow: [raw logits] -> mask -> [masked logits] -> argmax -> token
    \node[paneltitle] at (6.10, 4.95) {\large B.~Generation step};

    % --- (1) raw LM logits as 5 coloured bars (at top of the panel) ---
    \node[font=\scriptsize, color=black!70, anchor=base]
      at (6.10, 4.30) {LM logits};
    \foreach \i/\xx/\h in {0/5.00/0.55, 1/5.38/0.30, 2/5.76/0.45,
                            3/6.14/0.65, 4/6.52/0.25} {
      \draw[fill=c\i!75, draw=black!45, line width=0.4pt]
        (\xx, 3.55) rectangle (\xx + 0.28, 3.55 + \h);
    }
    \draw[axisline] (4.95, 3.55) -- (6.90, 3.55);

    % --- downward arrow labelled "mask V_{s+1}" ---
    \draw[flowarr] (6.10, 3.45) -- (6.10, 2.95)
       node[midway, right, font=\scriptsize, color=black!75]
       {mask $\mathcal V_{s+1}$};

    % --- (2) masked logits: only V_3 retained, others faded ---
    \foreach \i/\xx/\h in {0/5.00/0.55, 1/5.38/0.30, 2/5.76/0.45,
                            3/6.14/0.65, 4/6.52/0.25} {
      \draw[fill=c\i!75, draw=black!45, line width=0.4pt,
            opacity=0.18] (\xx, 2.20) rectangle (\xx + 0.28, 2.20 + \h);
    }
    % highlight V_3 (target partition, for example state s=2 -> s+1=3)
    \draw[fill=c3!75, draw=black!60, line width=0.7pt]
        (6.14, 2.20) rectangle (6.14 + 0.28, 2.20 + 0.65);
    \draw[axisline] (4.95, 2.20) -- (6.90, 2.20);

    % --- downward arrow labelled "arg max" ---
    \draw[flowarr] (6.10, 2.10) -- (6.10, 1.60)
       node[midway, right, font=\scriptsize, color=black!75]
       {$\arg\max$};

    % --- (3) output token chip in V_3 ---
    \node[bigpt, fill=c3!85] (toknext) at (6.10, 1.30) {};
    \node[font=\scriptsize, color=black!75, anchor=base]
      at (6.10, 0.85) {$t_{i+1} \in \mathcal V_{3}$};

    % =========== PANEL C: detection  (two stacked sequence plots) ===========
    %    Taller plots (y_step=0.30, height 1.20) for visual weight;
    %    symmetric around the mid-figure line so both plots feel balanced
    %    with the donut in Panel A.
    \def\xstep{0.55}
    \def\ystep{0.30}
    \def\plotL{9.90}
    \def\plotR{14.85}
    \def\plotw{4.95}
    \def\ploth{1.20}
    \def\pcen{0.60}   % y-offset from bot to centre of plot
    % title centred over the two stacked plots
    \node[paneltitle] at (\plotL + \plotw/2, 4.95)
      {\large C.~Detection by clockwork validity};

    % WATERMARKED plot at y in [3.20, 4.40]
    \def\wmbot{3.20}
    \def\wmtop{4.40}
    \node[font=\scriptsize\bfseries, color=green!45!black, anchor=base]
      at (\plotL + \plotw/2, \wmtop + 0.18)
      {Watermarked text:};

    \foreach \s in {0,1,2,3,4} {
      \draw[gridline] (\plotL, \wmbot + \s*\ystep) -- (\plotR, \wmbot + \s*\ystep);
    }
    \draw[axisline] (\plotL, \wmbot) -- (\plotR + 0.05, \wmbot);
    \draw[axisline] (\plotL, \wmbot) -- (\plotL, \wmtop + 0.05);
    \foreach \s in {0,1,2,3,4} {
      \node[yticklbl] at (\plotL - 0.05, \wmbot + \s*\ystep) {$\s$};
    }
    \node[axislabel, rotate=90, anchor=south]
      at (\plotL - 0.30, \wmbot + \pcen) {state};

    \draw[wmline]
      (\plotL + 0*\xstep, \wmbot + 0*\ystep)
      -- (\plotL + 1*\xstep, \wmbot + 1*\ystep)
      -- (\plotL + 2*\xstep, \wmbot + 2*\ystep)
      -- (\plotL + 3*\xstep, \wmbot + 3*\ystep)
      -- (\plotL + 4*\xstep, \wmbot + 4*\ystep)
      -- (\plotL + 5*\xstep, \wmbot + 0*\ystep)
      -- (\plotL + 6*\xstep, \wmbot + 1*\ystep)
      -- (\plotL + 7*\xstep, \wmbot + 2*\ystep)
      -- (\plotL + 8*\xstep, \wmbot + 3*\ystep)
      -- (\plotL + 9*\xstep, \wmbot + 4*\ystep);
    \foreach \i/\s in {0/0,1/1,2/2,3/3,4/4,5/0,6/1,7/2,8/3,9/4} {
      \node[pt, fill=c\s!85] at (\plotL + \i*\xstep, \wmbot + \s*\ystep) {};
    }
    \node[score, color=green!42!black]
      at (\plotR + 0.15, \wmbot + \pcen) {$\phi = 1$};

    % RANDOM plot at y in [1.00, 2.20]
    \def\rndbot{1.00}
    \def\rndtop{2.20}
    \node[font=\scriptsize\bfseries, color=red!60!black, anchor=base]
      at (\plotL + \plotw/2, \rndtop + 0.18)
      {Random text:};

    \foreach \s in {0,1,2,3,4} {
      \draw[gridline] (\plotL, \rndbot + \s*\ystep) -- (\plotR, \rndbot + \s*\ystep);
    }
    \draw[axisline] (\plotL, \rndbot) -- (\plotR + 0.05, \rndbot);
    \draw[axisline] (\plotL, \rndbot) -- (\plotL, \rndtop + 0.05);
    \foreach \s in {0,1,2,3,4} {
      \node[yticklbl] at (\plotL - 0.05, \rndbot + \s*\ystep) {$\s$};
    }
    \node[axislabel, rotate=90, anchor=south]
      at (\plotL - 0.30, \rndbot + \pcen) {state};

    \node[xticklbl] at (\plotL,                 \rndbot - 0.10) {$1$};
    \node[xticklbl] at (\plotL + 4.5*\xstep,    \rndbot - 0.10) {$\cdots$};
    \node[xticklbl] at (\plotL + 9*\xstep,      \rndbot - 0.10) {$n$};
    \node[axislabel, anchor=base]
      at (\plotL + \plotw/2, \rndbot - 0.55) {token index $i$};

    \draw[rndline]
      (\plotL + 0*\xstep, \rndbot + 3*\ystep)
      -- (\plotL + 1*\xstep, \rndbot + 0*\ystep)
      -- (\plotL + 2*\xstep, \rndbot + 2*\ystep)
      -- (\plotL + 3*\xstep, \rndbot + 4*\ystep)
      -- (\plotL + 4*\xstep, \rndbot + 1*\ystep)
      -- (\plotL + 5*\xstep, \rndbot + 3*\ystep)
      -- (\plotL + 6*\xstep, \rndbot + 0*\ystep)
      -- (\plotL + 7*\xstep, \rndbot + 4*\ystep)
      -- (\plotL + 8*\xstep, \rndbot + 2*\ystep)
      -- (\plotL + 9*\xstep, \rndbot + 1*\ystep);
    \foreach \i/\s in {0/3,1/0,2/2,3/4,4/1,5/3,6/0,7/4,8/2,9/1} {
      \node[pt, fill=c\s!85] at (\plotL + \i*\xstep, \rndbot + \s*\ystep) {};
    }
    \node[score, color=red!60!black]
      at (\plotR + 0.15, \rndbot + \pcen) {$\phi \approx 1/S$};

    % ----- short flow arrows between panels, centred in the gaps -----
    %   A->B: gap center x ~ 4.10; arrow 0.60 cm long
    %   B->C: gap center x ~ 8.40; arrow 0.80 cm long
    \draw[flowarr] (3.80, 2.30) -- (4.40, 2.30);
    \draw[flowarr] (8.00, 2.30) -- (8.80, 2.30)
       node[midway, above, font=\scriptsize, color=black!60]
       {append $t_{i+1}$};
  \end{tikzpicture}
  \caption{\textbf{How Markov Chain Lock works.}
    \textbf{(A)}~A keyed SHA-256 hash $\sigma_k$ partitions the
    vocabulary into $S{=}5$ disjoint sets $\mathcal V_0, \dots,
    \mathcal V_{S{-}1}$ ordered into a clockwork cycle
    $s \to s{+}1 \bmod S$.
    \textbf{(B)}~At a watermarked position with current state $s$,
    the language-model logits are masked to keep only the next
    partition $\mathcal V_{s+1}$, and the next token is the masked
    $\arg\max$.
    \textbf{(C)}~Detection re-derives every token's state from $k$
    in $O(n)$ hash operations: a watermarked sequence walks the
    chain by construction (top, $\phi = 1$), whereas unwatermarked
    text visits states uniformly and is valid only at rate $1/S$
    in expectation (bottom).}
  \label{fig:mcl_teaser}
\end{figure*}

We answer with \emph{Markov Chain Lock} (MCL) watermarking. MCL
sorts the vocabulary into $S$ colour-coded sets via a keyed
SHA-256 hash and forces a fraction $\rho$ of generated tokens to
walk a fixed cycle through those sets
(Figure~\ref{fig:mcl_teaser}). Detection re-derives every token's
colour from the key in $O(n)$ hash operations and counts how often
consecutive tokens follow the cycle. Two properties make the scheme
attractive for governance. First, detection runs without any
language-model access, so a third-party auditor needs neither the
provider's weights nor a proxy LM. Second, the detection bound
inverts in closed form to a calibration map: a regulator who fixes
a target false-positive rate $\alpha$ and a text-length floor
$n_{\min}$ reads off the minimum admissible state count
(Theorem~\ref{thm:detection}); a deployer can spend more budget for
margin but cannot cheat below the regulator's specification.

The choice of \emph{which} positions to watermark is left as a
design parameter and studied empirically. On Google's Gemma~3
270M model across 18 matched-budget cells
(three gates $\times$ three budgets $\times$ two domains), we find
that low-entropy gating, the policy implicit in
SWEET~\citep{lee2023sweet} and EWD~\citep{lu2024ewd}, is strictly
Pareto-dominated; that detection is gate-invariant at $\rho \geq
0.5$; and that the empirical post-attack scaling matches the
predicted $0.64$ on every cell. The operational dial is the budget
$\rho$, not the gate identity.

\textbf{Contributions.}
\begin{enumerate}[nosep, leftmargin=*]
  \item A discrete, hard-constraint watermarking primitive (MCL)
    with cryptographic state partitions, clockwork transitions, and
    $O(n)$ model-free detection (Section~\ref{sec:method}).
  \item A closed-form calibration formula
    (Theorem~\ref{thm:detection}) and a universal robustness
    threshold $\delta^\star = 1 - 1/\sqrt{2}$ under the midpoint
    decision convention (Theorem~\ref{thm:robustness}), strengthened
    to every oracle-blind adversary by self-healing
    (Theorem~\ref{thm:healing}); both identities extend to any
    $k$-regular doubly-stochastic transition topology
    (Proposition~\ref{prop:kregular}).
  \item Empirical evidence that low-entropy gating is dominated and
    that detection is gate-invariant at practical budgets
    (Section~\ref{sec:experiments}).
\end{enumerate}
